\subsubsection{多人性行为（3P / 群交 / 换偶）的安全与协商：当伴侣数大于二}

上一节讨论的是"一对一的偶遇"；当参与人数增加到三人及以上，风险不是线性增加，而是每多一个人就多一层交叉。本节适用于所有多人性情境——三人行（3P）、群交（group sex）、换偶（swinging）与多人派对，立场与全书一致：成年人自愿的多伴侣性行为本身不是道德问题，这里只谈协商与风险控制。

\textbf{第一层：伴侣间的协商先于一切}。与固定伴侣共同参与多人活动，最容易受伤的不是身体，而是"事后的信任"。事前必须明确三件事：
\begin{itemize}
\item \textbf{边界清单}：各自允许什么、不允许什么、哪些属于"看情况再说"的中间地带。边界可以随经验调整，但只能在事后复盘时调整——不能临场替对方"拓宽"；
\item \textbf{叫停信号}：约定任何一方都可以随时叫停，叫停不需要理由、不需要道歉，也不接受事后被追责"扫兴"；
\item \textbf{情绪预案}：多人场景最容易触发嫉妒、比较与被冷落感。提前约定事后的复盘时间；情绪反应是正常的，预先承认它比事后否认它更保护关系。
\end{itemize}

\textbf{第二层：屏障规则全面升级}。多伴侣情境下，"全程安全套"要升级为"一人一套、一事一套"：
\begin{itemize}
\item 阴茎更换性伴侣时\textbf{必须更换安全套}；口交同样使用安全套或口腔膜（dental dam）；
\item 玩具在不同的人之间使用前更换新的安全套套，或专人专用；手指接触建议使用指套或一次性手套；
\item 不要在口交接触不同人的生殖器后不经处理回到伴侣身上——口腔黏膜同样可以传播淋球菌、梅毒与 HPV；
\item 女性伴侣数增加时，细菌性阴道病与尿路感染的风险也随之上升；事后排尿、清水清洗是基础护理。
\end{itemize}

\textbf{第三层：检测频率相应提高}。暴露于多个新伴侣，意味着把检测（HIV、梅毒、淋病、衣原体）频率提高到每 3 个月一次，接种乙肝与 HPV 疫苗的价值也显著上升。同时要记住：HPV、疱疹与梅毒可以通过安全套未覆盖的皮肤接触传播——屏障降低了风险但没有归零，这正是检测不可替代的原因。如发生无保护的高危暴露，72 小时内按上一节的 PEP 流程处理。

\textbf{第四层：场景与清醒度}。多人派对是酒精与"派对药物"的高发情境，而同意能力与屏障坚持都依赖清醒——把"保持清醒"当作和"带够安全套"同级别的准备事项；饮品不离视线、警惕药物胁迫的规则同样适用（chemsex 的风险详见本书"性与药物"章节）。

最后强调：\textbf{同意必须覆盖每一个环节}。照片与录像在多人场景中是高发纠纷源——未经同意拍摄或传播他人性行为影像涉嫌违法；如果现场有人提议拍摄，任何一方都可以否决。
